\par \SourceLeaf{462}


und ferner
\[
\begin{gathered}
48.\quad{}
(1-2x^3+2x^{12}- .\,. )+i(1-2x+2x^4- .\,. )\\[8pt]
\quad=1-x\,.\,1-x^3\,.\,1-x^5\mathinner{.\,.\,.} \\[8pt]\quad{}1+i\,.\,1+x^3\,.\,1-x^6\,.\,1+x^9\mathinner{.\,.\,.}\\[8pt]
\qquad\\[8pt]\quad{}\times1-ix\,.\,1-ixx\,.\,\\[8pt]\quad{}1+ix^4\,.\,\\[8pt]\quad{}1+ix^5\,.\,1-ix^7\,.\,1-ix^8\mathinner{.\,.\,.}
\end{gathered}
\]
\[
\begin{gathered}
49.\quad{}
(1+2x^3+2x^{12}+ .\,. )-i(1+2x+2x^4+ .\,. )\\[8pt]
\quad=1+x\,.\,1+x^3\,.\,1+x^5\mathinner{.\,.\,.} \\[8pt]\quad{}1-i\,.\,1-x^3\,.\,1-x^6\,.\,1-x^9\mathinner{.\,.\,.}\\[8pt]
\qquad\\[8pt]\quad{}\times\\[8pt]\quad{}1-ix\,.\,1+ixx\,.\,\\[8pt]\quad{}1-ix^4\,.\,1+ix^5\mathinner{.\,.\,.}
\end{gathered}
\]
\[
\begin{gathered}
50.\quad{}
(1+2x^3+2x^{12}+ .\,. )+i(1+2x+2x^4+ .\,. )\\[8pt]
\quad=1+x\,.\,1+x^3\,.\,1+x^5\mathinner{.\,.\,.} \\[8pt]\quad{}1+i\,.\,1-x^3\,.\,1-x^6\,.\,1-x^9\mathinner{.\,.\,.}\\[8pt]
\qquad\\[8pt]\quad{}\times\\[8pt]\quad{}1+ix\,.\,1-ixx\,.\,\\[8pt]\quad{}1+ix^4\,.\,1-ix^5\mathinner{.\,.\,.}
\end{gathered}
\]

Also aus der Multiplication von 47 und 48
\[
\begin{gathered}
51.\quad{}
(1-2x^3+2x^{12}-\mathinner{.\,.\,.})^2+(1-2x+2x^4-\mathinner{.\,.\,.})^2\\[8pt]
\quad=2(1-x)^2(1-x^3)^2(1-x^5)^2\mathinner{.\,.\,.}\\[8pt]\quad{}(1+x^3)^2(1-x^6)^2(1+x^9)^2\mathinner{.\,.\,.}\\[8pt]
\qquad\times1+xx\,.\,1+x^4\,.\,1+x^8\,.\,1+x^{10}\mathinner{.\,.\,.}\\[8pt]
\quad=2(1-x)^2(1-x^3)^2(1-x^5)^2\mathinner{.\,.\,.}\\[8pt]\quad{}(1+x^3)^2 1+x^9)^2\mathinner{.\,.\,.}\\[8pt]
\qquad\times(1-x^6)^3(1-x^{12})^2\\[8pt]\quad{}\mathinner{.\,.\,.} 1+xx\,.\,1+x^4\,.\,1+x^6\,.\,1+x^8\mathinner{.\,.\,.}\\[8pt]
\quad=2\frac{[x]^2[x^6]^4}{[xx]^2[x^3]^2[x^{12}]^2}\,.\,\frac{[x^6]^3}{[x^{12}]}\,.\,\frac{[x^4]}{[xx]}
\\[8pt]\quad{}=\frac{2[x]^2[x^4][x^6]^7}{[x^2]^3[x^3]^2[x^{12}]^3}
\end{gathered}
\]

und aus der Multiplication von 49 und 50
\[
\begin{gathered}
52.\quad{}
(1+2x^3+2x^{12}+ .\,. )^2+(1+2x+2x^4+\mathinner{.\,.\,.})^2\\[8pt]
\quad=2(1+x)^2(1+x^3)^2(1+x^5)^2\mathinner{.\,.\,.}\\[8pt]\quad{}(1-x^3)^2(1-x^6)^2(1-x^9)^2\mathinner{.\,.\,.}\\[8pt]
\qquad\times1+xx\,.\,1+x^4\,.\,1+x^8\,.\,1+x^{10}\mathinner{.\,.\,.}\\[8pt]
\quad=2\frac{[xx]^4}{[x]^2[x^4]^2}[x^3]^2\frac{[x^4][x^6]}{[xx][x^{12}]}
\\[8pt]\quad{}=\frac{2[xx]^3[x^3]^2[x^6]}{[x]^2[x^4][x^{12}]}
\end{gathered}
\]

und der Quotient
\[
\begin{gathered}
53.\quad{}
\frac{(1-2x^3+2x^{12}- .\,. )^2+(1-2x+2x^4- .\,. )^2}
{(1+2x^3+2x^{12}+ .\,. )^2+(1+2x+2x^4+ .\,. )^2}\\[8pt]
\quad=\Bigl(\frac{1-x}{1+x}\Bigr)^2\Bigl(\frac{1-x^3}{1+x^3}\Bigr)^2\Bigl(\frac{1-x^5}{1+x^5}\Bigr)^2\mathinner{.\,.\,.}
\\[8pt]\quad{}\Bigl(\frac{1+x^3}{1-x^3}\Bigr)^2\Bigl(\frac{1+x^9}{1-x^9}\Bigr)^2\mathinner{.\,.\,.}\\[8pt]
\quad=\Bigl(\frac{1-x}{1+x}\Bigr)^2\Bigl(\frac{1-x^5}{1+x^5}\Bigr)^2\Bigl(\frac{1-x^7}{1+x^7}\Bigr)^2\mathinner{.\,.\,.}
\\[8pt]\quad{}=\frac{[x]^4[x^4]^2[x^6]^6}{[x^2]^6[x^3]^4[x^{12}]^2}
\end{gathered}
\]
\SourceLeaf{463}


und das Product
\[
\begin{gathered}
54.\quad{}
\{(1-2x^3+ .\,. )^2+(1-2x+ .\,. )^2\}\\[8pt]\quad{}\ \{(1+2x^3+ .\,. )^2-(1+2x+ .\,. )^2\}\\[8pt]
\quad=4\frac{[x^6]^8}{[x^{12}]^4}=4(1-2x^6+2x^{24}- .\,. )^2
\end{gathered}
\]

Aus Formel 23 folgt
\[
\begin{gathered}55.\quad{}
x+x^9+x^{25}+ .\,. \\[8pt]\quad{}=x\,.\,1+x^8\,.\,1+x^{16}\,.\,1+x^{24}\mathinner{.\,.\,.}\\[8pt]\quad{}(1-x^{16}-x^{32}+x^{80}+x^{112}- .\,. )\end{gathered}
\]
\[
\frac32\ \text{Exponent}=\square-1
\]

Aus Formel 26 folgt
\[
\begin{gathered}56.\quad{}
x^3+x^{27}+x^{75}+ .\,. \\[8pt]\quad{}=x\,.\,1+x^8\,.\,1+x^{16}\,.\,1+x^{24}\mathinner{.\,.\,.}\\[8pt]\quad{}(x^2-x^{10}-x^{42}+x^{66}+x^{130}- .\,. )\end{gathered}
\]
\[
\begin{gathered}\text{oder}\quad x^{\frac13}\,.\,1+x^8\,.\,1+x^{16}\,.\,1+x^{24}\mathinner{.\,.\,.}\\[8pt]\quad{}=A,\quad x^2=t^3\quad\text{gesetzt}\end{gathered}
\]
\[
\begin{gathered}55.\quad{}
x+x^9+x^{25}+ .\,. \\[8pt]\quad{}=A(t-t^{25}-t^{49}+t^{121}+t^{169}-\mathinner{.\,.\,.})\end{gathered}
\]
\[
\begin{gathered}56.\quad{}
x^3+x^{27}+x^{75}+ .\,. \\[8pt]\quad{}=A(t^4-t^{16}-t^{64}+t^{100}+t^{196}-\mathinner{.\,.\,.})\end{gathered}
\]

Nun folgt aus der Factorenzerlegung in 24 sehr leicht, wenn man statt \(x\), \(it\)
statt \(y\), \(i\) setzt
\[
\begin{gathered}
it-it^4+it^{16}-it^{25}-it^{49}+ .\,.\\[8pt]
\quad=it\,.\,1-t^3\,.\,1+t^{15}\,.\,1+t^{21}\,.\,1-t^{33}\mathinner{.\,.\,.}\\[8pt]\quad{}1+t^{18}\,.\,1-t^{36}\,.\,1+t^{54}\mathinner{.\,.\,.}
\end{gathered}
\]

Also aus 55---56
\[
\begin{gathered}
57.\quad{}
(x+x^9+x^{25}+ .\,. )-(x^3+x^{27}+x^{75}+ .\,. )\\[8pt]
\quad=x\,.\,1-xx\,.\,1+x^{10}\,.\,1+x^{14}\,.\,1-x^{22}\,.\,1-x^{26}\mathinner{.\,.\,.}\\[8pt]
\qquad\times1+x^{12}\,.\,1-x^{24}\,.\,1+x^{36}\mathinner{.\,.\,.}\\[8pt]\quad{}1+x^8\,.\,1+x^{16}\,.\,1+x^{24}\mathinner{.\,.\,.}
\end{gathered}
\]

Und so sehr leicht
\[
\begin{gathered}
58.\quad{}
(x+x^9+x^{25}+ .\,. )+(x^3+x^{27}+x^{75}+ .\,. )\\[8pt]
\quad=x\,.\,1+xx\,.\,1-x^{10}\,.\,1-x^{14}\,.\,1+x^{22}\,.\,1+x^{26}\mathinner{.\,.\,.}\\[8pt]
\qquad\times1+x^{12}\,.\,1-x^{24}\,.\,1+x^{36}\mathinner{.\,.\,.}\\[8pt]\quad{}1+x^8\,.\,1+x^{16}\,.\,1+x^{24}\mathinner{.\,.\,.}
\end{gathered}
\]

Also durch Multiplication
\SourceLeaf{464}

\[
\begin{gathered}
59.\quad{}
(x+x^9+x^{25}+ .\,. )^2-(x^3+x^{27}+x^{75}+ .\,. )^2\\[8pt]
\quad=xx\,.\,1-x^4\,.\,1-x^{20}\,.\,1-x^{28}\,.\,1-x^{44}\,.\,1-x^{52}\mathinner{.\,.\,.}\\[8pt]
\qquad\times(1+x^{12})^2(1-x^{24})^2(1+x^{36})^2\mathinner{.\,.\,.}\\[8pt]\quad{}(1+x^8)^2(1+x^{16})^2\mathinner{.\,.\,.}\\[8pt]
\quad=xx\frac{[x^4][x^{16}]^2[x^{24}]^7}{[x^8]^3[x^{12}]^3[x^{48}]^2}
\end{gathered}
\]

Eben so folgt aus der Factorenzerlegung in 24, wenn man statt \(x\), \(t\) und
statt \(y\), \(-i\) setzt
\[
\begin{gathered}t+it^4-it^{16}-t^{25}-t^{49}\mathinner{.\,.\,.}\\[8pt]\quad{}=t\,.\,1+it^3\,.\,1-it^{15}\,.\,1+it^{21}\,.\,1-it^{33}\mathinner{.\,.\,.}[t^{18}]\end{gathered}
\]

also
\[
\begin{gathered}
60.\quad{}
(x+x^9+x^{25}+\mathinner{.\,.})+i(x^3+x^{27}+x^{75}+\mathinner{.\,.})\\[8pt]
=x\,.\,1+ixx\,.\,1-ix^{10}\,.\,1+ix^{14}\,.\,1-ix^{22}\mathinner{.\,.\,.}\\[8pt]\quad{}1+x^8\,.\,1+x^{16}\mathinner{.\,.\,.}\\[8pt]\quad{}1-x^{12}\,.\,1-x^{24}\mathinner{.\,.\,.}
\end{gathered}
\]

und eben so
\[
\begin{gathered}
61.\quad{}
(x+x^9+x^{25}+\mathinner{.\,.})-i(x^3+x^{27}+x^{75}+\mathinner{.\,.})\\[8pt]
=x\,.\,1-ixx\,.\,1+ix^{10}\,.\,1-ix^{14}\,.\,1+ix^{22}\mathinner{.\,.\,.}\\[8pt]\quad{}1+x^8\,.\,1+x^{16}\mathinner{.\,.\,.}\\[8pt]\quad{}1-x^{12}\,.\,1-x^{24}\mathinner{.\,.\,.}
\end{gathered}
\]

Also durch Multiplication
\[
\begin{gathered}
62.\quad{}
(x+x^9+x^{25}+\mathinner{.\,.})^2+(x^3+x^{27}+ .\,. )^2\\[8pt]
\quad=xx\frac{[x^8]^2}{[x^4][x^{16}]}\,.\,\frac{[x^{12}][x^{48}]}{[x^{24}]^2}\,.\,\frac{[x^{16}]^2}{[x^8]^2}\,.\,[x^{12}]^2
\\[8pt]\quad{}=xx\frac{[x^{12}]^3[x^{16}][x^{48}]}{[x^4][x^{24}]^2}
\end{gathered}
\]

Also Product von 59 und 62
\[
\begin{gathered}63.\quad{}
(x+x^9+x^{25}+\mathinner{.\,.})^4-(x^3+x^{27}+x^{75}+\mathinner{.\,.})^4
\\[8pt]\quad{}=x^4\frac{[x^{16}]^3[x^{24}]^5}{[x^8]^3[x^{48}]}\end{gathered}
\]

Durch Multiplication von 62 und 44 folgt
\[
\begin{gathered}
64.\quad{}
\{(x^{\frac12}+x^{\frac92}+x^{\frac{25}2}+ .\,. )^2+(x^{\frac32}+x^{\frac{27}2}+ .\,. )^2\}\\[8pt]
\quad\times\{(1+2x^2+2x^8+ .\,. )^2-(1+2x^6+2x^{24}+ .\,. )^2\}\\[8pt]
\quad=4xx\frac{[x^4]^3[x^{12}][x^{24}]^2}{[x^2][x^6][x^8]^2}\,.\,x\frac{[x^6]^3[x^8][x^{24}]}{[x^2][x^{12}]^2}
\\[8pt]\quad{}=4x^3\frac{[x^4]^3[x^6]^2[x^{24}]^3}{[x^2]^2[x^{12}][x^8]}
\end{gathered}
\]

Durch Multiplication von 59 und 44

\SourceLeaf{465}

\[
\begin{gathered}
65.\quad{}
\{(x^{\frac12}+x^{\frac92}+x^{\frac{25}2}+\mathinner{.\,.})^2
 -(x^{\frac32}+x^{\frac{27}2}+\mathinner{.\,.})^2\}\\[8pt]
\quad\times\{(1+2x^2+2x^8+\mathinner{.\,.})^2
 -(1+2x^6+2x^{24}+\mathinner{.\,.})^2\}\\[8pt]
\qquad=4x^3\frac{[xx][x^8]^2[x^{12}]^7}{[x^4]^3[x^6]^3[x^{24}]^2}
 \cdot\frac{[x^4]^3[x^{12}][x^{24}]^2}{[x^2][x^6][x^8]^2}\\[8pt]
\qquad=4x^3\frac{[x^{12}]^8}{[x^6]^4}
 \\[8pt]\quad{}=4(x^{\frac32}+x^{\frac{27}2}+x^{\frac{75}2}+\mathinner{.\,.})^4
\end{gathered}
\]
Man hat ferner
\[
\begin{gathered}66.\quad{}
[x]=\frac{a^{\frac1{12}}b^{\frac13}(aa-bb)^{\frac1{24}}}
 {2^{\frac13}x^{\frac1{24}}\mathord{\surd}h}
\\[8pt]\quad{}\qquad
\frac ah=\frac{[x^2]^{10}}{[x]^4[x^4]^4}\end{gathered}
\]
\[
\begin{gathered}67.\quad{}
[xx]=\frac{a^{\frac16}b^{\frac16}(aa-bb)^{\frac1{12}}}
 {2^{\frac23}x^{\frac1{12}}\mathord{\surd}h}
\\[8pt]\quad{}\qquad
\frac bh=\frac{[x]^4}{[xx]^2}\end{gathered}
\]
\[
\begin{gathered}68.\quad{}
[x^4]=\frac{a^{\frac1{12}}b^{\frac1{12}}(aa-bb)^{\frac16}}
 {2^{\frac43}x^{\frac16}\mathord{\surd}h}
\\[8pt]\quad{}\qquad
\frac{aa-bb}{hh}=16x\frac{[x^4]^8}{[xx]^4}\end{gathered}
\]

\textit{Für die Fünf-Theilung}\quad
$\left\{\begin{array}{cc}a&b\\A&B\end{array}\right\}$
\[
\begin{gathered}
\Bigl(\frac{a-A}{5A-a}\Bigr)^4=\frac{AA-BB}{aa-bb}\cdot\frac{BB}{bb}\\[8pt]
\Bigl(\frac{B-b}{5B-b}\Bigr)^4=\frac{AA-BB}{aa-bb}\cdot\frac{AA}{aa}\\[8pt]
\frac{a-A}{B-b}=\frac{5B-b}{5A-a}=\mathord{{}^{4}\!\surd}{\frac{aB}{bA}}
\end{gathered}
\]
\[
\begin{gathered}
(a-A)(5A-a)^5=256Aabb(aa-bb)\\[8pt]
(B-b)(5B-b)^5=256Baab(aa-bb)
\end{gathered}
\]
\[
\begin{gathered}
(a-A)^5(5A-a)=256AaBB(AA-BB)\\[8pt]
(B-b)^5(5B-b)=256AAbB(AA-BB)
\end{gathered}
\]
\[
\begin{gathered}
a-A=\frac{4x[xx]^2[x^5][x^{20}]}{[x][x^4]}
 \\[8pt]\quad{}=2^{\frac43}\frac{a^{\frac16}A^{\frac16}B^{\frac5{12}}(AA-BB)^{\frac5{24}}}
 {b^{\frac1{12}}(aa-bb)^{\frac1{24}}}\\[1ex]
5A-a=4\frac{[x][x^4][x^{10}]^2}{[x^5][x^{20}]}
 \\[8pt]\quad{}=2^{\frac43}\frac{a^{\frac16}A^{\frac16}b^{\frac5{12}}(aa-bb)^{\frac5{24}}}
 {B^{\frac1{12}}(AA-BB)^{\frac1{24}}}
\end{gathered}
\]

\SourceLeaf{466}

Zu der Theorie der Fünftheilung gehören folgende Theoreme. Wir bezeichnen
$1+xy\,.\,1+x^3y\,.\,1+x^5y\mathinner{.\,.\,.}1+\frac xy\,.\,1+\frac{x^3}y\,.\,1+\frac{x^5}y\mathinner{.\,.\,.}$
durch $(x,y)$, so ist
\[
\begin{gathered}
[69]\quad{}
(x,\alpha y).(x,\tfrac ya)
\\[8pt]\quad{}=P\{1+xx(yy+\tfrac1{yy})+x^8(y^4+\tfrac1{y^4})+\mathinner{.\,.}\}\\[8pt]
\quad+Q\{y+\tfrac1y+x^4(y^3+\tfrac1{y^3})+x^{12}(y^5+\tfrac1{y^5})+\mathinner{.\,.}\}
\end{gathered}
\]
wo $P$, $Q$ von $y$ unabhängig.

Also
\[
\begin{gathered}(x,\alpha i)(x,\tfrac ai)=(xx,\alpha\alpha)
 \\[8pt]\quad{}=P(1-2xx+2x^8-\mathinner{.\,.})=\frac{[xx]^2}{[x^4]}P\end{gathered}
\]
oder
\[
P=(xx,\alpha\alpha)\frac{[x^4]}{[xx]^2}
\]
Ferner für $y=-\alpha x$
\[
\begin{gathered}
(x,-\alpha\alpha x).(x,-x)=0\\[8pt]
=P(1+\tfrac1{aa}+\alpha\alpha x^4+\tfrac{x^4}{a^4}+\mathinner{.\,.})
 \\[8pt]\quad{}-Q(\alpha x+\tfrac1{ax}+\tfrac x{a^3}+\alpha^3x^7+\tfrac{x^7}{a^5}+\mathinner{.\,.\,.})
\end{gathered}
\]
d.\ i.
\[
\begin{gathered}
P\{\alpha+\tfrac1a+x^4(\alpha^3+\tfrac1{a^3})+x^{12}(\alpha^5+\tfrac1{a^5})+\mathinner{.\,.}\}\\[8pt]
\quad=\frac Qx\{1+xx(\alpha\alpha+\tfrac1{aa})+x^8(\alpha^4+\tfrac1{a^4})+\mathinner{.\,.}\}
\end{gathered}
\]
Nun ist
\[
P=\frac1{[xx]^2}\{1+xx(\alpha\alpha+\tfrac1{aa})+\mathinner{.\,.}\}
\]
Also
\[
Q=\frac x{[xx]}\{(\alpha+\tfrac1a)+x^4(\alpha^3+\tfrac1{a^3})+\mathinner{.\,.\,.}\}
\]
und unser Theorem
\[
\begin{gathered}
70.\quad{}
(x,\alpha y).(x,\tfrac ya)\\[8pt]
=\frac1{[xx]^2}\left\{\begin{gathered}
\phantom{{}+x}\{1+xx(\alpha\alpha+\tfrac1{aa})+x^8(\alpha^4+\tfrac1{a^4})+\mathinner{.\,.}\}\\[8pt]
\quad\times\{1+xx(yy+\tfrac1{yy})+x^8(y^4+\tfrac1{y^4})+\mathinner{.\,.}\}\\[1ex]
+x\{\alpha+\tfrac1a+x^4(\alpha^3+\tfrac1{a^3})+x^{12}(\alpha^5+\tfrac1{a^5})+\mathinner{.\,.}\}\\[8pt]
\quad\times\{y+\tfrac1y+x^4(y^3+\tfrac1{y^3})+x^{12}(y^5+\tfrac1{y^5})+\mathinner{.\,.}\}
\end{gathered}\right\}
\end{gathered}
\]


\SourceLeaf{467}

oder
\[
\begin{gathered}(x,\alpha y).(x,\tfrac ya)
 =\frac{[x^4]^2}{[xx]^2}\{(xx,\alpha\alpha).(xx,yy)
 \\[8pt]\quad{}+x\alpha y(xx,\alpha\alpha xx).(xx,xxyy)\}\end{gathered}
\]
(Man kann auch leicht die Reihe, wodurch $P$ multiplicirt ist, $=0$ machen,
durch $y=ix$)
\[
71.\quad{}
(x,y)+(x,\tfrac xy)\mathord{{}^{4}\!\surd}{\frac x{yy}}
 =(x^{\frac14},y^{\frac12})\frac{[x^{\frac12}]}{[xx]}
\]
\[
\begin{gathered}(x,\alpha x)=\frac1a(x,\tfrac xa),
\\[8pt]\quad{}\qquad
(x,\alpha xx)=\frac1{ax}(x,\alpha)\end{gathered}
\]
Den Satz 70 kann man auch so enonciren
\[
\begin{gathered}72.\quad{}
(x,\alpha).(x,\beta)=\frac{[x^4]^2}{[xx]^2}
 \{(xx,\alpha\beta).(xx,\tfrac a\beta)
 \\[8pt]\quad{}+x\alpha(xx,\alpha\beta xx).(xx,\tfrac{axx}\beta)\}\end{gathered}
\]
Hieraus folgt
\[
\begin{gathered}(x,\tfrac xa).(x,\tfrac x\beta)=\frac{[x^4]^2}{[xx]^2}
 \{(xx,\tfrac{xx}{a\beta}).(xx,\tfrac a\beta)
 \\[8pt]\quad{}+\alpha(xx,\alpha\beta).(xx,\tfrac{axx}\beta)\}\end{gathered}
\]
hieraus ferner
\[
\begin{gathered}73.\quad{}
(x,\alpha).(x,\beta)+(x,\tfrac xa).(x,\tfrac x\beta)\mathord{\surd}{\frac x{a\beta}}
 \\[8pt]\quad{}=\frac{[x]^2}{[xx]^2}.(x^{\frac12},\mathord{\surd}{\alpha\beta}).(x^{\frac12},\mathord{\surd}{\tfrac\beta a})\end{gathered}
\]
Nun ist
\[
\begin{gathered}
\frac{1+2x+2x^4+\mathinner{.\,.\,.}}{1+2x^5+2x^{20}+\mathinner{.\,.\,.}}
\\[8pt]\quad{}=\frac{(ix^{\frac52},ix^{\frac32}).(ix^{\frac52},-ix^{\frac12})}
 {(ix^{\frac52},-ix^{\frac32}).(ix^{\frac52},ix^{\frac12})}\\[1ex]
\\[8pt]\quad{}=\frac{(-x^5,xx)(-x^5,-x)+x(-x^5,-x^3)(-x^5,x^4)}
 {(-x^5,xx)(-x^5,-x)-x(-x^5,-x^3)(-x^5,x^4)}
\end{gathered}
\]
Woraus der erste zu beweisende Satz von selbst folgt. Ebenso ist
\[
\begin{gathered}
\frac{1+2x+2x^4+\mathinner{.\,.\,.}}{1+2x^5+2x^{20}+\mathinner{.\,.\,.}}\\[8pt]
\quad=\dfrac{\begin{gathered}1-\varepsilon x\,.\,1-\varepsilon\varepsilon x\,.\,1-\varepsilon^3x\,.\,1-\varepsilon^4x\,.\,\\[6pt]1+\varepsilon xx\,.\,1+\varepsilon\varepsilon xx\,.\,1+\varepsilon^3xx\,.\,1+\varepsilon^4xx\mathinner{.\,.\,.}\end{gathered}}
 {\begin{gathered}1+\varepsilon x\,.\,1+\varepsilon\varepsilon x\,.\,1+\varepsilon^3x\,.\,1+\varepsilon^4x\,.\,\\[6pt]1-\varepsilon xx\,.\,1-\varepsilon\varepsilon xx\,.\,1-\varepsilon^3xx\,.\,1-\varepsilon^4xx\mathinner{.\,.\,.}\end{gathered}}\\[1ex]
\quad=\frac{1-\varepsilon^4\,.\,1-\varepsilon^3}{1+\varepsilon^4\,.\,1+\varepsilon}
 \cdot\frac{(ix^{\frac12},i\varepsilon x^{\frac12}).(ix^{\frac12},i\varepsilon\varepsilon x^{\frac12})}
 {(ix^{\frac12},-i\varepsilon x^{\frac12}).(ix^{\frac12},i\varepsilon\varepsilon x^{\frac12})}\\[1ex]
\quad=\frac{\varepsilon\varepsilon-\varepsilon^3\,.\,\varepsilon-\varepsilon^4}
 {\varepsilon\varepsilon+\varepsilon^3\,.\,\varepsilon+\varepsilon^4}
 \cdot\frac{(-x,-\varepsilon^3x).(-x,\varepsilon)-\varepsilon x(-x,\varepsilon^3xx)(-x,-\varepsilon x)}
 {(-x,-\varepsilon^3x).(-x,\varepsilon)+\varepsilon x(-x,\varepsilon^3xx)(-x,-\varepsilon x)}\\[1ex]
\quad=\frac{-\varepsilon+\varepsilon\varepsilon+\varepsilon^3-\varepsilon^4}
 {+\varepsilon+\varepsilon\varepsilon+\varepsilon^3+\varepsilon^4}
 \cdot\frac{(x,\varepsilon^3x).(x,-\varepsilon)+\varepsilon^3(x,-\varepsilon^3)(x,\varepsilon x)}
 {(x,\varepsilon^3x).(x,-\varepsilon)-\varepsilon^3(x,-\varepsilon^3)(x,\varepsilon x)}
\end{gathered}
\]
Woraus der zweite zu beweisende Satz von selbst folgt.

\SourceLeaf{468}

\[
\begin{gathered}
(x,x)=2(1+xx)^2(1+x^4)^2\mathinner{.\,.\,.}
       =2\frac{[x^4]^2}{[xx]^2}\\[8pt]
(x,1)=(1+x)^2\ (1+x^3)^2\mathinner{.\,.\,.}
       =\frac{[xx]^4}{[x]^2[x^4]^2}
\end{gathered}
\]
Man hat also
\[
\begin{gathered}(x,\alpha)^2=\frac{[x^4]^2}{[xx]^2}
\Bigl\{(xx,\alpha\alpha)\frac{[x^4]^4}{[xx]^2[x^8]^2}
\\[8pt]\quad{}+2x\alpha(xx,\alpha\alpha xx)\frac{[x^8]^2}{[x^4]^2}\Bigr\}\end{gathered}
\]

Durch die Entwickelung von $(x,y)^3$ erhält man
\[
\begin{gathered}P\Bigl\{1+x^3\Bigl(y^3+\frac{1}{y^3}\Bigr)+.\,.\Bigr\}
\\[8pt]\quad{}+Q\Bigl\{x^{\frac13}\Bigl(y+\frac1y\Bigr)
\\[8pt]\quad{}+x^{\frac43}\Bigl(yy+\frac1{yy}\Bigr)
+x^{\frac{16}3}\Bigl(y^4+\frac1{y^4}\Bigr)+.\,.\Bigr\}\end{gathered}
\]
also für $y=-\varepsilon x$
\[
\begin{gathered}(1-\varepsilon\varepsilon)^3\frac{[x^6]^3}{[x^2]^3}
\\[8pt]\quad{}=Qx^{-\frac23}\{\varepsilon-\varepsilon\varepsilon
-(\varepsilon-\varepsilon\varepsilon)xx-\mathinner{.\,.\,.}\}\end{gathered}
\]
oder
\[
\begin{gathered}3x^{\frac23}.\frac{[x^6]^3}{[x^2]^3}
\\[8pt]\quad{}=Q(1-xx-x^4+x^{10}+x^{14}-\mathinner{.\,.\,.})=[xx]Q\end{gathered}
\]
also
\[
Q=3x^{\frac23}.\frac{[x^6]^3}{[xx]^4}
\]

Durch Entwickelung von $(x,y).(x,-y)^2$ erhält man
\[
\begin{gathered}P'\Bigl\{1+x^3\Bigl(y^3+\frac1{y^3}\Bigr)+\mathinner{.\,.\,.}\Bigr\}
\\[8pt]\quad{}+Q'\Bigl\{x^{\frac13}\Bigl(y+\frac1y\Bigr)
+x^{\frac43}\Bigl(yy+\frac1{yy}\Bigr)+.\,.\Bigr\}\end{gathered}
\]
\[
\begin{gathered}(1-\varepsilon\varepsilon)(1+\varepsilon\varepsilon)^2
\frac{[x^6][x^8][x^{12}]^3}{[xx][x^4]^3[x^{24}]}
\\[8pt]\quad{}=Q'x^{-\frac23}(\varepsilon-\varepsilon\varepsilon)[xx]\end{gathered}
\]
\[
Q'=-x^{\frac23}.\frac{[x^6][x^8][x^{12}]^3}{[xx]^2[x^4]^3[x^{24}]}
\]

Man folgert hieraus leicht, dass man setzen darf
\[
T(x,y)^3+U(x,y)(x,-y)^2=(x^3,y^3)
\]
so dass $T$ und $U$ Functionen von $x$. Um sie zu bestimmen setzen wir


\SourceLeaf{469}

\[
\begin{gathered}1)\qquad y=x\qquad\text{so wird}\qquad
T(x,x)^3=(x^3,x^3)\\[8pt]\quad{}\quad\text{oder}\quad
T=\tfrac14\frac{[x^{12}]^2[xx]^6}{[x^6]^2[x^4]^6}\end{gathered}
\]
\[
2)\qquad y=-\varepsilon x\quad\text{so wird}\quad
T(x,-\varepsilon x)^2=U(x,+\varepsilon x)^2
\]
oder
\[
T\frac{(1-\varepsilon\varepsilon)^2[x^6]^2}{[xx]^2}
=\frac{[x^{12}]^2[xx]^2}{[x^6]^2[x^4]^2}U(1+\varepsilon\varepsilon)^2
\]
oder
\[
T=-\tfrac13 U\frac{[xx]^4[x^{12}]^2}{[x^6]^4[x^4]^2}
\]
oder
\[
U=-\tfrac34\frac{[x^6]^2[xx]^2}{[x^4]^4}
\]
folglich
\[
\begin{gathered}\frac{T\Bigl[\dfrac{(x,-y)}{(x,y)}\Bigr]^3
+U\dfrac{(x,-y)}{(x,y)}}
{T+U\Bigl[\dfrac{(x,-y)}{(x,y)}\Bigr]^2}
=\frac{(x^3,-y^3)}{(x^3,y^3)}\end{gathered}
\]
also
\[
\begin{gathered}\frac{\dfrac{[xx]^4}{[x^4]^2}
\Bigl[\dfrac{(x,-y)}{(x,y)}\Bigr]^3
-3\dfrac{[x^6]^4}{[x^{12}]^2}\dfrac{(x,-y)}{(x,y)}}
{\dfrac{[xx]^4}{[x^4]^2}
-3\dfrac{[x^6]^4}{[x^{12}]^2}
\Bigl[\dfrac{(x,-y)}{(x,y)}\Bigr]^2}
=\frac{(x^3,-y^3)}{(x^3,y^3)}\end{gathered}
\]
\begin{center}\rule{.23\textwidth}{.3pt}\end{center}

